\section{Discussion}

The results confirm the core analytical prediction (delay destabilizes) and reveal a surprising empirical finding (learning buffers rather than amplifies). This section discusses three qualitative insights that go beyond the numerical results: why the ODE-to-simulation gap exists and what it teaches us about the mechanism, what the architecture hierarchy reveals about the nature of delay vulnerability, and what practical implications follow for institutional design.

\subsection{Why the quantitative gap between theory and simulation matters (Q1)}

The ODE predicts a sharp bifurcation at $\Delta_c$; the simulation shows a gradual dose-response curve. This gap is not a failure of the theory---it is informative about the mechanism. Three stabilizing buffers are absent from the mean-field model and present in the simulation: stochastic exploration ($\epsilon$-greedy noise damps incipient oscillations), the three-action space (agents can retreat to moderate behavior rather than oscillating between extremes), and spatial heterogeneity (network structure disperses the synchronized oscillations required for large-amplitude instability). These are not defects of the simulation; they are what the mean-field theory trades away for a closed-form answer.

The qualitative lesson is that delay-induced instability is robust to realistic complications, but the sharp threshold of the ODE becomes a gradual erosion of stability in a heterogeneous population. Institutions operating near the theoretical boundary should not expect a clean phase transition; they should expect slowly worsening oscillatory behavior that may be difficult to distinguish from normal fluctuations until runaway is already underway.

\subsection{What the architecture hierarchy reveals about the mechanism (Q3)}

The ordering (reactive agents most fragile, Q-learning agents partially resilient, fixed-policy agents immune) points to a specific causal mechanism rather than a generic correlation between complexity and instability.

The mechanism is a feedback loop: when the alarm drops (due to institutional processing lag), reactive agents immediately escalate. Punishment arrives too late; by then the alarm spike triggers the next oscillation. The critical feature is not that reactive agents are ``simple'' but that they are \emph{memoryless}: each decision is based entirely on the current (stale) signal, with no integration over past experience. Q-learning agents break this loop because their Q-values carry forward the memory of past punishment. The ``soft brake'' is not a design feature---it is an emergent consequence of temporal-difference learning applied to a delayed-feedback environment.

This distinction has a practical implication: the risk factor for delay vulnerability is not whether agents are adaptive, but whether they integrate information over time. Any decision process that reacts only to current signals (whether a simple threshold, a rule-based policy, or a sophisticated model with no memory state) will be vulnerable. Institutions seeking stability should therefore prioritize reducing processing delay over restricting agent adaptiveness. The problem is the lag, not the learning.

\subsection{Policy levers: response sharpness and adaptive governance (Q2, Q4)}

The theory predicts that gradual institutional responses (low $k$) expand the stability margin. The simulation confirms this, but with an important qualification: the protective effect of gradual response is conditional on agent architecture. Fixed-policy agents are stable regardless of sharpness, because they do not close the feedback loop. For adaptive agents, reducing sharpness helps but cannot eliminate delay vulnerability---it only raises the threshold. The policy implication is that institutions face a genuine dilemma: sharp responses are decisive but fragile; gradual responses are resilient but slow to act. Every parameter trades one vulnerability for another.

The second lever is the regulator itself. The RL regulator experiment is exploratory, and the specific numbers should be interpreted cautiously. The qualitative pattern, however, is suggestive: the regulator appears to convert catastrophic runaway into bounded oscillations rather than restoring full stability. This mode conversion---disaster into nuisance---may be the realistic ceiling for any controller that shares the institution's information delay. A regulator that observes the same stale signal as the repression mechanism cannot anticipate the population's response; it can only react to the consequences of its own delayed actions, introducing a secondary feedback loop that sustains limit cycles while preventing unbounded growth. We report this as a pattern worth investigating, not a confirmed finding: 50 seeds is enough to see a pattern but not enough to bet on it.

\subsection{Limitations}

This model is deliberately simple, and the simplifications cost something. The Q-learning model is deliberately minimal and does not capture sophisticated strategic reasoning, communication, or coordination among agents. Real adaptive agents in social or technical systems may exhibit richer behavioral repertoires. The three-action space is a simplification; continuous action spaces might produce different stability characteristics. The regime classifier uses fixed thresholds, and the precise quantitative rates depend on these choices. We verified that the central ordering (fixed $\leq$ Q-learning $\leq$ reactive in delay sensitivity) is preserved across runaway thresholds from $R_{\max}>0.30$ through $R_{\max}>0.50$, though the absolute rates shift substantially (e.g., reactive at delay${}=20$ ranges from 100\% at threshold 0.30 to 64\% at threshold 0.50). A convergence diagnostic confirms that Q-learning agents reach stable behavior by step 100--200: the mean radical fraction and its standard deviation change by less than 0.001 between the last two 100-step windows at all tested delays, indicating that 500 steps is sufficient for the Q-values to equilibrate.

The simulation uses finite populations ($N=240$) on specific graph realizations; larger populations or different graph families might shift the quantitative boundaries. The architecture hierarchy itself, however, is robust to the network parameterization: re-running the crossed delay~$\times$~architecture design on a denser modular graph (4 communities, $p_{\mathrm{in}}=0.15$, $p_{\mathrm{out}}=0.02$) reproduces the same ordering---fixed agents immune (0\% at all delays), reactive agents catastrophic (100\% by delay~8), Q-learning agents partially resilient (77\% at delay~20)---with the denser network somewhat more unstable (reactive reaches 27\% runaway even at delay~0), consistent with its stronger influence coupling.

The reactive baseline uses a single threshold heuristic; other reactive architectures (imitation dynamics, best-response with noise) might show different levels of delay sensitivity, though the qualitative finding (that memoryless reactivity is more fragile than learning) should hold for any policy that lacks temporal integration. The RL regulator's Q-update bootstraps toward a next-state computed from the current (undelayed) alarm rather than the next delayed observation, introducing a minor state-transition inconsistency; since Experiment~6 is presented as exploratory evidence, this does not affect the paper's confirmed findings. Finally, the theory assumes a single aggregate alarm signal; systems with multiple observation channels or local feedback might exhibit different instability structures. Reality, of course, does not present its instabilities one at a time.

\subsection{Relationship to prior work}

Our analytical results connect directly to the Hopf bifurcation analysis of \citet{wesson2016hopf} for two-strategy delayed replicator dynamics, extending their framework to include an asymmetric institutional response function rather than symmetric frequency-dependent payoffs. The discrete-time instability amplification we observe in Experiment~2 echoes the finding of \citet{alboszta2004stability} that delay can destabilize the ESS in discrete replicator dynamics; \citet{iijima2012delayed} documents the complementary regime in which stability is instead preserved, underscoring that the destabilizing effect of delay is model-dependent. The network simulation builds on the evolutionary-dynamics-on-graphs tradition \citep{lieberman2005evolutionary, ohtsuki2006simple, szabo2007evolutionary}, adding delayed institutional feedback as a novel destabilizing mechanism distinct from the structural effects studied in that literature. The interaction between learning and delay connects to recent work on reward delays in multi-agent reinforcement learning \citep{zhang2023marl}, though our focus on population-level regime transitions rather than individual convergence is novel.

The companion paper extends the present mechanism to noisy selective control on modular networks, addressing the question of how imperfect observation and targeted governance alter stability conditions. That extension is necessary because real institutions rarely observe true system state directly, and the consequences of classification errors depend strongly on network position, a consideration we deliberately set aside here so that delay gets a fair hearing on its own.
